El flujo eléctrico de un campo eléctrico $\mathbf E$ a través de una superficie orientada $S$ se define por
\[
\Phi_{\rm e}=\iint_S\mathbf E\cdot\mathbf n\,dS,
\]
donde $\mathbf n$ es la normal unitaria elegida. Una esfera maciza de radio $R>0$, centrada en el origen, tiene carga distribuida uniformemente con densidad volumétrica $\rho_0$ (carga por unidad de volumen). En su interior, el campo está dado por
\[
\mathbf E(x,y,z)=\frac{\rho_0}{3\varepsilon_0}\langle x,y,z\rangle,
\qquad x^2+y^2+z^2<R^2,
\]
donde $\varepsilon_0>0$ es la permitividad eléctrica del vacío, una constante que relaciona la carga con el campo eléctrico. Para $0<a<R$, calcule mediante el teorema de la divergencia el flujo eléctrico exterior por $S_a:x^2+y^2+z^2=a^2$.